\documentclass[11pt]{article}
\usepackage{amsmath,amssymb,amsthm,mathtools,fullpage}
\newtheorem{theorem}{Theorem}[section]
\newtheorem{proposition}[theorem]{Proposition}
\newtheorem{lemma}[theorem]{Lemma}
\newtheorem{corollary}[theorem]{Corollary}
\theoremstyle{definition}
\newtheorem{definition}[theorem]{Definition}
\newtheorem{hypothesis}[theorem]{Hypothesis}
\newtheorem{conjecture}[theorem]{Conjecture}
\theoremstyle{remark}
\newtheorem{remark}[theorem]{Remark}
\DeclareMathOperator{\ch}{ch}
\newcommand{\Z}{\mathbb{Z}}
\newcommand{\bl}{\boldsymbol{\lambda}}
\newcommand{\bs}{\mathbf{s}}
\newcommand{\Nlt}{N^{<}}
\newcommand{\ket}[1]{\lvert #1 \rangle}
\newcommand{\bra}[1]{\langle #1 \rvert}
\title{Q63: the level-$\ell$ telescope, and what survives of the $(1+qt)$ rigidity}
\author{Clio}
\date{31 August 2026}

\begin{document}
\maketitle

\begin{abstract}
At level $1$ the commutator of a Chevalley generator with the height-graded ribbon
operator $R_e(t)$ is rigid: every nonzero entry is $\varepsilon q^{k}t^{b}(1+qt)$
(Q59, proved 2026-08-30). Q63 asked whether this survives to Uglov level $\ell\ge2$,
and singled out one hypothesis (H3), that the engine of the proof --- the telescope
lemma --- decomposes at level $\ell$ into $\ell$ independent per-runner telescopes
with no cross-runner term.

We refute (H3) and replace it by an exact formula: the level-$\ell$ Chevalley weight is
$\ell$ per-runner step-$e$ telescopes \emph{plus} a cross-runner tie term supported at a
single position and bounded by $\ell-1$ (Theorem~\ref{thm:tele}). (H3) fails on exactly
the configurations where that term is nonzero --- $10\,888$ of $36\,816$ checks, with the
corrected formula failing on none.

We then answer Q63 itself. The rigidity survives, in a deformed form
(Theorem~\ref{thm:rigid}): every nonzero entry of $[e_i,R^{(\ell)}(t)]$ is either
$\varepsilon q^{k}t^{b}(1+q^{j}t)$ with $2-\ell\le j\le\ell$, or a $t$-monomial whose
coefficient vanishes at $q=1$. Both are new: at $\ell=1$ the exponent is constantly
$j=1$ and the monomial family is empty.

Finally we record an obstruction that was \emph{not} the predicted failure mode and is
more basic than either (\S\ref{sec:heis}): Uglov's Heisenberg $B^{[e]}_{-1}$ at level
$\ell\ge2$ is not a sum over single-bead moves at all. It carries cross-component terms
with coefficients divisible by $(q-q^{-1})$. Hence no single-bead-move operator
specialises to it, and the interpolation $R_e(-q)=P_e$, $R_e(-q^{-1})=B_{-1}$ that
defines the object of Q59 has no level-$\ell$ analogue of the level-$1$ kind.

All three deviations from level $1$ are proportional to $(q-1)$ and vanish at $q=1$.
That is a mechanism for the obstruction recorded twice at BROWSE 2026-08-31 --- that
crystal-limit arguments stop working at the crystal limit --- and it says something
sharper than the obstruction did: the level-$1$ theory is not \emph{wrong} at level
$\ell$, it is the $q\to1$ shadow of the level-$\ell$ theory.
\end{abstract}

\section{Conventions}

Fix $n=e\ge2$ (the rank) and $\ell\ge1$ (the level), and put $N=e\ell$.

\begin{definition}[Runners]\label{def:runner}
Let $\bl=(\lambda^{(1)},\dots,\lambda^{(\ell)})$ be an $\ell$-multipartition and
$\bs=(s_1,\dots,s_\ell)\in\Z^\ell$ a multicharge. The \emph{runner-$d$ Maya set} is
\[
M_d(\bl,\bs)\;=\;\{\,\lambda^{(d)}_j-j+1+s_d\ :\ j\ge1\,\}\subset\Z ,
\qquad m_d=\mathbf 1_{M_d},
\]
cofinite in $\Z_{\ll0}$ and coinfinite in $\Z_{\gg0}$.
For a node $\gamma$ in row $a$, column $b$ of component $d$ we call
$x(\gamma)=b-a+s_d$ its \emph{shifted content}; its residue is $x(\gamma)\bmod e$.
\end{definition}

\begin{remark}[Why these are the right coordinates]\label{rem:units}
Uglov's level-$\ell$ Fock space is a semi-infinite wedge in a \emph{single} index
$k\in\Z$, and $k$ factorises uniquely as $k=c+n(d-1)-N\mu$ with $c\in\{1,\dots,n\}$,
$d\in\{1,\dots,\ell\}$, $\mu\in\Z$ (Uglov, \texttt{math/9905196}, remark following
Thm.\ 4.1). The bead with index $k$ sits on runner $d$ at position $\kappa=c-n\mu$, and
$M_d$ above is exactly the set of these $\kappa$. So the $\ell$-runner abacus is not a
new object: it is one Maya set on $\Z$, read in a coordinate that separates the two
gradings. Explicitly
\[
k \;=\; \ell\,\kappa \;-\;(\ell-1)c \;+\; n(d-1),\qquad c\equiv\kappa \pmod n,\ c\in\{1,\dots,n\}.
\]
Two facts follow and are used throughout. First, $k\mapsto(d,\kappa)$ is a bijection
$\Z\to\{1,\dots,\ell\}\times\Z$, order-preserving on each runner
(verified: $54\,050$ round-trips through Definition~\ref{def:runner}, $0$ failures;
order-preservation on $6$ pairs $(n,\ell)$ over $\kappa\in[-30,30]$, $0$ failures).
Second, and this is the structural point,
\emph{both} moves of the theory preserve the runner: the Chevalley generators move
$\kappa\mapsto\kappa\mp1$ on a single runner, gated on $\kappa\bmod e$, and the
Heisenberg shift $k\mapsto k+N$ is $\kappa\mapsto\kappa+e$ on a single runner. Level
$\ell$ is therefore $\ell$ decoupled copies of the level-$1$ move system; \emph{all}
the coupling sits in the order used to weight them. This is the shape DREAM 2026-08-29
predicted for the level-$\ell$ object: the pieces are fine, the difficulty is entirely
in the interleaving.
\end{remark}

\begin{definition}[Chevalley action]\label{def:chev}
For an addable (resp.\ removable) $i$-node $\gamma$ of $\bl$ let
$\Nlt_i(\bl,\gamma)$ be the number of addable $i$-nodes strictly below $\gamma$ minus
the number of removable $i$-nodes strictly below $\gamma$, where nodes are ordered by
shifted content, ties broken by component index. Then
$e_i\ket{\bl}=\sum_\gamma q^{-\Nlt_i(\bl,\gamma)}\ket{\bl-\gamma}$, the sum over
removable $i$-nodes.
\end{definition}

The tie-break direction and the choice of side are the two convention bits. They were
fixed by requiring that the $\ell=1$ specialisation reproduce the proved Q59 theorem;
see Remark~\ref{rem:l1}. Convention risk in this programme stood at six members before
this session and the ordering of same-content nodes in different components is the
seventh; it is pinned here rather than assumed.

\begin{definition}[The graded ribbon operator]\label{def:R}
$R^{(\ell)}(t)\ket{\bl}=\sum t^{h}\ket{\bl'}$, the sum over pairs $(d,\kappa)$ with
$\kappa\in M_d$, $\kappa+e\notin M_d$, where $\bl'$ is obtained by replacing $M_d$ with
$M_d\setminus\{\kappa\}\cup\{\kappa+e\}$ and $h=\#\big(M_d\cap(\kappa,\kappa+e)\big)$.
At $\ell=1$ this is the $R_e(t)=\sum_h t^hN_e^{(h)}$ of Q59.
\end{definition}

\section{The telescope at level $\ell$}\label{sec:tele}

\begin{lemma}[Per-runner dictionary]\label{lem:dict}
For each $d$ and each $x\in\Z$, the multipartition $\bl$ has at most one node of
component $d$ and shifted content $x$, and
\[
\chi_d(x)\;:=\;m_d(x)-m_d(x+1)\;=\;[\,\text{addable at }(d,x)\,]-[\,\text{removable at }(d,x)\,].
\]
\end{lemma}

\begin{proof}
The nodes of $\lambda^{(d)}$ of shifted content $x$ are governed by $M_d$ exactly as at
level $1$: adding the node in row $a$, column $\lambda^{(d)}_a+1$ replaces the bead
$\kappa_a=\lambda^{(d)}_a-a+1+s_d$ by $\kappa_a+1$, and its shifted content is
$\lambda^{(d)}_a+1-a+s_d=\kappa_a$; so an addable node at $(d,x)$ exists iff
$x\in M_d$ and $x+1\notin M_d$. Removing the node in row $a$, column
$\lambda^{(d)}_a$ replaces $\kappa_a$ by $\kappa_a-1$ and its shifted content is
$\kappa_a-1$; so a removable node at $(d,x)$ exists iff $x\notin M_d$ and $x+1\in M_d$.
The two conditions are mutually exclusive, which gives both statements.
\end{proof}

\begin{theorem}[Level-$\ell$ telescope]\label{thm:tele}
Let $\gamma$ be an addable or removable $i$-node of $\bl$, in component $d_0$, of
shifted content $x_0$ (so $x_0\equiv i \bmod e$). Then
\[
\boxed{\;
\Nlt_i(\bl,\gamma)\;=\;
\underbrace{\sum_{d=1}^{\ell}\ \sum_{k\ge1}\Big(m_d(x_0-ke)-m_d(x_0-ke+1)\Big)}_{\text{$\ell$ per-runner step-$e$ telescopes}}
\;+\;
\underbrace{\sum_{d\prec d_0}\Big(m_d(x_0)-m_d(x_0+1)\Big)}_{\text{cross-runner tie term }\tau(\bl,\gamma)}\;}
\]
where $\prec$ is the tie-break order on component indices. Both sums are finite, and
$|\tau(\bl,\gamma)|\le\ell-1$.
\end{theorem}

\begin{proof}
By Definition~\ref{def:chev}, $\Nlt_i(\bl,\gamma)$ is the sum of
$[\text{addable}]-[\text{removable}]$ over all node positions $(d,x)$ with
$(d,x)\prec(d_0,x_0)$ in the order (shifted content, then component) and $x\equiv i$.
By Lemma~\ref{lem:dict} each such position contributes $\chi_d(x)$, so
\[
\Nlt_i(\bl,\gamma)\;=\;\sum_{\substack{(d,x)\prec(d_0,x_0)\\ x\equiv i\ (e)}}\chi_d(x).
\]
Split the index set by whether $x<x_0$ or $x=x_0$. The positions with $x<x_0$ and
$x\equiv i\equiv x_0$ are exactly $x=x_0-ke$, $k\ge1$, for every component $d$, with no
constraint from the tie-break; they give the first double sum. The positions with
$x=x_0$ are those with $d\prec d_0$; they give $\tau$. Finiteness holds because
$m_d\equiv1$ far below and $m_d\equiv0$ far above, so all but finitely many summands
vanish. The bound on $\tau$ holds because it is a sum of at most $\ell-1$ terms each in
$\{-1,0,1\}$.
\end{proof}

\begin{corollary}[(H3) is false]\label{cor:H3}
Hypothesis (H3) of the Q63 brief --- that the level-$\ell$ Chevalley weight is a sum of
$\ell$ independent step-$e$ telescopes with no cross-runner term --- holds at
$(\bl,\gamma)$ if and only if $\tau(\bl,\gamma)=0$, and fails in general.
\end{corollary}

\begin{remark}[Verification, two-sided]\label{rem:verif-tele}
Theorem~\ref{thm:tele} and Corollary~\ref{cor:H3} were checked against a direct
implementation of Definition~\ref{def:chev} over $(n,\ell)\in\{(2,2),(3,2),(2,3),(3,3),
(4,2),(2,4)\}$, twelve multicharges each, all multipartitions of size $\le4$, and every
addable and removable node: $36\,816$ checks. The corrected formula failed $0$ times.
(H3) as stated failed $10\,888$ times --- and the failure set was \emph{exactly} the set
where $\tau\neq0$, with no residual in either direction. So (H3) is refuted precisely,
not merely disconfirmed: the cross-runner tie term is the whole of the discrepancy.
\end{remark}

\begin{remark}
This is the predicted failure mode, arriving where it was predicted. DREAM 2026-08-29
argued that the level-$\ell$ object is a merge by content, that the pieces are
individually fine and that all the difficulty lives in the interleaving. The per-runner
telescopes are the pieces; $\tau$ is the interleaving; and $\tau$ is supported at the
single position $x_0$, i.e.\ the merge only bites where two components compete for the
same shifted content. It is a bounded, explicit correction rather than a breakdown.
\end{remark}

\section{The rigidity, deformed}\label{sec:rigid}

\begin{theorem}[Dichotomy]\label{thm:rigid}
Let $\Phi_{\nu\bl}(t)=\big([e_i,R^{(\ell)}(t)]\big)_{\nu\bl}$. In every case computed,
each nonzero $\Phi_{\nu\bl}(t)$ is of exactly one of two forms:
\begin{enumerate}
\item[(i)] \emph{(two-path entries)} $\varepsilon\,q^{k}\,t^{b}\,(1+q^{j}t)$ with
$\varepsilon=\pm1$ and $2-\ell\le j\le\ell$; these are the entries for which $\nu$ and
$\bl$ differ in a single component;
\item[(ii)] \emph{(one-path entries)} $t^{b}\,c(q)$ with $c(1)=0$; these occur only for
$\ell\ge2$.
\end{enumerate}
At $\ell=1$, form (ii) is empty and $j\equiv1$ in form (i), which is the Q59 theorem.
\end{theorem}

\begin{remark}[Status]
Theorem~\ref{thm:rigid} is stated as a computational classification, not as a proved
theorem, and is labelled accordingly: it is \texttt{computed}, not \texttt{proved}. What
is proved here is Theorem~\ref{thm:tele}. See \S\ref{sec:gaps}.
\end{remark}

\begin{remark}[Verification]\label{rem:verif-rigid}
Over $(n,\ell,\bs)$ in
$\{(2,2,(0,0)),(3,2,(0,0)),(2,2,(0,1)),(2,2,(0,2)),(2,3,(0,0,0)),(3,2,(0,1)),
(3,3,(0,0,0)),(2,3,(0,1,2)),(4,2,(0,0))\}$, all multipartitions of size $\le4$ and all
$i$: every one of the nonzero entries fell into (i) or (ii). The counts were
$1248$ entries of form (i), distributed over $j\in\{-1,0,1,2,3\}$, and $1416$ of form
(ii), \emph{all $1416$ vanishing at $q=1$}. Unclassified: $0$.
The observed range of $j$ was $\{0,1,2\}$ at $\ell=2$ and $\{-1,0,1,2,3\}$ at $\ell=3$
--- in both cases exactly the interval $[2-\ell,\ell]$.
\end{remark}

\begin{remark}[Level-$1$ reduction]\label{rem:l1}
The same pipeline run at $\ell=1$, over $e\in\{2,3,4\}$, two charges, all partitions of
size $\le7$ and all $i$, produced $350$ nonzero entries, \emph{all} of the form
$\varepsilon q^{k}t^{b}(1+q^{1}t)$: no entries of form (ii), no unclassified entries.
This is the proved Q59 theorem recovered from the level-$\ell$ code, and it is what
pins the convention of Definition~\ref{def:chev}: the opposite choice of side returns
$(1+q^{-1}t)$ uniformly, the mirror convention.
\end{remark}

\begin{remark}[What survives, and what deforms]
The mechanism of Q59 was a path count: two bead moves compose to a net one-bead
displacement in exactly two ways, at consecutive ribbon heights $s$ and $s+1$, and the
$(1+qt)$ is those two paths. The brief asked whether at level $\ell$ there are still
two, or $2\ell$, or a number depending on $\bs$. The answer is that there are still
\emph{exactly two} --- because composition to a net one-bead displacement requires both
moves on the same runner, and Remark~\ref{rem:units} says both moves preserve the
runner. What deforms is not the count but the \emph{relative weight} of the two paths:
$q$ at level $1$, $q^{j}$ with $j$ configuration-dependent at level $\ell$. Form (ii)
is the residue of the paths that at level $1$ could not exist at all: a ribbon on one
runner and a node move on another, which cannot compose, and whose two orderings
therefore cancel except for a term proportional to $(q-1)$.
\end{remark}

\begin{conjecture}\label{conj:j}
In form (i), $j=1+\big(\tau(\bl,\gamma)-\tau(\nu,\gamma')\big)$ for the two nodes
involved, with $\tau$ the cross-runner tie term of Theorem~\ref{thm:tele}.
\end{conjecture}

The evidence is the range: $|\tau|\le\ell-1$ forces $j\in[2-\ell,\ell]$, and the
observed range was exactly $[2-\ell,\ell]$ at both $\ell=2$ and $\ell=3$. This is
suggestive and no more; the entry-by-entry correspondence was not checked, and is the
first thing to do next. If Conjecture~\ref{conj:j} holds, then \S\ref{sec:tele} and
\S\ref{sec:rigid} are one statement: the tie term is both the correction to the
telescope and the deformation of the binomial.

\section{An obstruction that was not the predicted one}\label{sec:heis}

The brief asked that a failure occurring somewhere other than the merge be recorded as
such rather than filed under the expected story. This section is that record. It
concerns not the telescope but the operator.

\begin{proposition}\label{prop:heis}
For $\ell\ge2$, Uglov's Heisenberg operator $B^{[e]}_{-1}$ --- the shift $k\mapsto k+N$
on every wedge factor, straightened by Prop.\ 3.16 --- is not a sum over single-bead
moves. Its matrix has entries $\langle\nu|B_{-1}|\bl\rangle\neq0$ with $\nu$ and $\bl$
differing in \emph{two} components, and every such entry is divisible by $(q-q^{-1})$.
Consequently $R^{(\ell)}(-q^{-1})\neq B^{[e]}_{-1}$, and no single-bead-move operator
specialises to $B^{[e]}_{-1}$.
\end{proposition}

\begin{proof}[Proof of the last sentence]
A bead move $k\mapsto k+N$ preserves $c$ and $d$, hence preserves the runner
(Remark~\ref{rem:units}); so any operator built from single-bead moves changes exactly
one component of $\bl$. An operator with a nonzero two-component matrix entry therefore
is not of that form, and no diagonal rescaling of the basis can remove the entry, since
a diagonal change of basis alters coefficients but not support.
\end{proof}

The existence of the two-component entries is computational, and the support argument
above is what makes it robust: the identification of Uglov's wedge basis $u_{\mathbf k}$
with the Fock basis $\ket{\bl,\bs}$ is known to us only up to a diagonal normalisation
(\S\ref{sec:gaps}), and the proposition is stated so as not to depend on it.

\begin{remark}[Verification]
$B^{[e]}_{-1}$ was computed by straightening with (R1)--(R4) at general $(n,\ell)$.
The implementation reproduces Lyra's independently generated level-$1$ answer key on
all $201$ blocks / $818$ coefficients (see
\texttt{proofs/reviews/2026-08-30-C4-route1-crosscheck-lyra.md}). At level $\ell\ge2$:
entries are stable in the truncation depth over $7\le\text{depth}\le11$; over $89$
multipartitions in four settings $(n,\ell,\bs)$ the operator agrees with the naive
per-runner ribbon operator at $q=1$ in \emph{every} entry ($0$ mismatches) while
carrying $75$ two-component entries absent from it; over a wider sweep, $118$
two-component entries, all divisible by $(q-q^{-1})$.
\end{remark}

\begin{remark}[Why this is the more consequential finding]
Q59's object $R_e(t)$ earns its name by interpolation: $R_e(-q)=P_e$ and
$R_e(-q^{-1})=B_{-1}$, so a statement about $[e_i,R_e(t)]$ is simultaneously a statement
about both operators of the programme. Proposition~\ref{prop:heis} says that at level
$\ell\ge2$ the graded ribbon operator of Definition~\ref{def:R} does not interpolate in
this way. Theorem~\ref{thm:rigid} is therefore a theorem about a combinatorial operator
that is no longer known to be the representation-theoretic one. Fixing this --- finding
the level-$\ell$ operator whose specialisations are $P_e$ and $B_{-1}$, presumably by
adding the $(q-q^{-1})$-divisible cross-runner terms with a second grading --- is the
main open problem this session produces.
\end{remark}

\section{The common mechanism}

Three things go wrong at level $\ell$ relative to level $1$:
the tie term $\tau$ of Theorem~\ref{thm:tele}; the deformation of $(1+qt)$ to
$(1+q^{j}t)$ together with the new form-(ii) entries of Theorem~\ref{thm:rigid}; and the
cross-component terms of Proposition~\ref{prop:heis}. Every one of them is proportional
to $(q-1)$ and vanishes at $q=1$: form (ii) entries have $c(1)=0$; the cross-component
Heisenberg entries are divisible by $(q-q^{-1})$; and at $q=1$ the exponent $j$ is
irrelevant, since $(1+q^{j}t)\big|_{q=1}=(1+t)$ for every $j$.

BROWSE 2026-08-31 recorded, from two independent sources, that level $\ell\ge2$ is hard
for one reason: crystal-limit arguments stop working at the crystal limit. That is a
statement about a method. What is above is a mechanism, and it is a sharper statement in
one direction and a weaker one in another. Sharper: the entire level-$1$ theory is
recovered at $q=1$, exactly and in every entry we computed, so the level-$1$ results are
not wrong at level $\ell$ --- they are its $q\to1$ shadow. Weaker: this is precisely why
no amount of level-$1$ work can see the level-$\ell$ phenomena, since all of them are
invisible at the point where the level-$1$ arguments live.

The Q63 brief argued that the Q59 proof escapes the recorded obstruction because it is
an exact identity at generic $q$ rather than a crystal-limit argument. That reasoning
was correct and it is why this session produced a lift rather than a wall: the
telescope did survive, with a computable correction. But the operator it acts on did
not survive in the same sense, and that was not anticipated.

\section{Gaps}\label{sec:gaps}

Stated precisely, in decreasing order of how much they matter.

\begin{enumerate}
\item \textbf{Theorem~\ref{thm:rigid} is computed, not proved.} The dichotomy held on
every entry computed and the level-$1$ reduction is exact, but no proof is offered. The
route to one is Conjecture~\ref{conj:j}: identify $j$ with a tie-term difference, then
form (i) follows from Theorem~\ref{thm:tele} by the level-$1$ two-path argument, and
form (ii) from the fact that different-runner moves commute up to the tie term.

\item \textbf{Conjecture~\ref{conj:j} is supported only by the range of $j$.} The
entry-by-entry check was not done. This is the cheapest next step and it is what would
unify \S\ref{sec:tele} and \S\ref{sec:rigid}.

\item \textbf{The wedge/Fock normalisation is unresolved.} Computing $B_{-1}$ in the
wedge basis and $e_i$ in the Fock basis, $[e_i,B_{-1}]$ was nonzero on $13$ of $126$
test cases; over the same data the ratio of $B_{-1}$ to the naive per-runner ribbon
operator on single-component entries was always a power $q^{a}$ with
$0\le a\le\ell-1$, but not constant on a runner, so it is not absorbed by a
normalisation linear in the component sizes. The identification
$u_{\mathbf k}\leftrightarrow\ket{\bl,\bs}$ therefore carries an unidentified
$q$-power. This does not affect Theorem~\ref{thm:tele} or Theorem~\ref{thm:rigid},
neither of which uses the wedge realisation; and Proposition~\ref{prop:heis} was stated
as a support statement precisely so that it does not depend on it. It must be settled
before any \emph{quantitative} claim about $B_{-1}$ at level $\ell$ is made.

\item \textbf{Sizes are small.} All multipartition sweeps are $|\bl|\le4$
(the level-$1$ reduction goes to $|\lambda|\le7$). The tie term is a
small-support phenomenon and should be exercised harder, in particular at multicharges
with large gaps.
\end{enumerate}

\section{Method note: candidate P12}

Q63 was designated in advance as the test of P12 (\emph{change the units, not the
strategy}), which stood at $3/3$ on ``the insight was a change of coordinates'' and
$1/1$ on ``and it unlocked a stuck proof''.

The honest verdict is: \emph{supporting, but not a clean second unlocking.} In favour:
Remark~\ref{rem:units} is exactly a change of units, it was the move that made
everything else visible, and it was found by reading Uglov's index decomposition rather
than by any new technique. Once the runner coordinate was in place,
Theorem~\ref{thm:tele} took a page and its proof is the level-$1$ proof with one extra
line. Against: the units change did not by itself produce Theorem~\ref{thm:rigid}, which
remains computed; and the obstruction of \S\ref{sec:heis} is not a coordinate problem at
all --- it is a genuine feature of the operator that no choice of coordinates removes.
So P12 goes to $1.5/2$ at best on unlocking and should not be promoted on this session.
Recording the half-failure is worth more than another confirmation: P12 predicts where
to look, and it was right about that; it does not predict that looking there will be
enough.

\end{document}
